\documentclass[addpoints]{exam}
% \documentclass[addpoints,12pt]{exam}

\usepackage[slovene]{babel}
\usepackage{amsfonts}
\usepackage{amssymb}
\usepackage{amsmath}
\usepackage{float}
\usepackage{graphicx}
\usepackage{enumitem}
\usepackage{color}
\usepackage{eurosym}

\usepackage{pgfpages}
% \usepackage{tikz}
\usepackage{wrapfig}
\usepackage{pgfkeys}
\usepackage{pgfplots}
\usepackage{xcolor}
\usepackage{tkz-euclide}
\usepackage{xfp}




%Komentar naslednje vrstice glede na verzijo testa -- rešitve ali prostor za odgovore:
% \printanswers

% \linespread{2}


\pointsinrightmargin
\marginbonuspointname{}


\renewcommand{\solutiontitle}{\noindent\textbf{Rešitve in točkovanje:}\par\noindent}

\colorgrids
\definecolor{GridColor}{gray}{0.7}
\setlength{\gridsize}{7mm}

\extrawidth{5mm}
\setlength{\rightpointsmargin}{1.5cm}

\begin{document}

\runningfooter{}{}{}

\begin{center}
    \fbox{\fbox{\parbox{15cm}{\centering
    Popravljanje in pridobivanje ocen -- drugo pisno ocenjevanje znanja \\ 2. letnik -- test G \\ 23. januar 2026 \\ (Kotne funkcije; Vektorji)}}}
\end{center}

\vspace{0.1in}
\makebox[\textwidth]{Ime in priimek, oddelek:\enspace\hrulefill}


\begin{center}
    \hqword{Naloga}
    \hpword{Možne točke}
    % \chbpword{Dodatne točke}
    \hsword{Dosežene točke}
    \htword{Skupaj}  
    % \combinedpointtable[h][questions]
    \gradetable[h][questions]

\end{center}

\vspace{0.1in}


\begin{questions}

    % 1. vprašanje

    \question[4]
    Poenostavite izraz. 
    $$\displaystyle \left(\left(\tan{x}\cos{x}\right)^{-2}+\cos^{-2}{x}\right)\sin^2{x} $$

    \begin{solution}[6.5cm]
       
    \end{solution}


    % 2. vprašanje

    \question[5]
    Izračunajte natančno vrednost izraza. 
    $$\displaystyle \dfrac{\left(\cos{135^\circ}+\sin{225^\circ}\right)^2}{\tan{300^\circ}-\tan{120^\circ}-\sin{270^\circ}} $$

    \begin{solution}[6cm]
       
    \end{solution}


\newpage
    % 3. vprašanje

    \question
    V pravilnem šestkotniku $ABCDEF$ točka $M$ tretjini stranico $BC$, točka $N$ razdeli stranico $ED$ v razmerju $|EN|:|ND|=2:3$.
    Z baznima vektorjema $\overrightarrow{a}=\overrightarrow{AB}$ in $\overrightarrow{b}=\overrightarrow{AF}$ izrazite vektorje: 
    \begin{parts}
        \part[3]
        $\overrightarrow{MN}$, 

    \begin{solution}[5cm]

    \end{solution}
        
        \part[2] 
        $\overrightarrow{NB}-\overrightarrow{FB}+\overrightarrow{FA}-\overrightarrow{EA}$.
    \end{parts}

    \begin{solution}[4cm]

    \end{solution}

    

    
    % 4. vprašanje

    \question[5]
    V trikotniku $\triangle ABC$ z oglišči $A(4,5,-6)$, $B(1,7,0)$ in $C(4,0,-3)$ izračunajte oddaljenost težišča trikotnika $\triangle ABC$ od razpolovišča stranice $AB$.
    % \begin{parts}

    %     \part[2]
    %     dolžino težiščnice na stranico $AC=b$, 

    % \begin{solution}[6cm]

    % \end{solution}

    %     \part[3] 
    %     koordinate težišča trikotnika $\triangle ABC$.

    % \begin{solution}[2cm]

    % \end{solution}        
    
    % \part[5]
    %     skalarni produkt $\overrightarrow{AB}\cdot\overrightarrow{AC}$ in kot $\angle BAC=\alpha$, 

    %     \begin{solution}[5cm]

    % \end{solution}




    % \end{parts}







    \newpage
    % 5. vprašanje

    \question
    Daljica $AB$ je določena s krajiščema $A(5,-2,7)$ in $B(1,5,11)$.
    
    \begin{parts}
        \part[4]
        Izračunajte skalarni produkt krajevnega vektorja točke $A$ in vektorja  $2\overrightarrow{AB}$.
       
           \begin{solution}[6cm]

    \end{solution}

    %     \part[3]    
    %     Določite enotski vektor v smeri krajevnega vektorja točke $S$, ki je razpolovišče daljice $AB$.

    % \begin{solution}[5cm]

    % \end{solution}

        \part[3]
        Določite $x$, da bo vektor $\overrightarrow{v}=(2x,7,x+8)$ pravokoten na vektor $-\overrightarrow{AB}$.

            \begin{solution}[6cm]

    \end{solution}

        \part[3]
        Določite $y$ in $z$, da bo vektor $\overrightarrow{u}=(z+1,10,y)$ vzporeden krajevnemu vektorju točke $A$.
    \end{parts}

\begin{solution}[4cm]

\end{solution}


\newpage

    % 6. vprašanje



    \question[4]
    Dan je trikotnik $\triangle ABC$ s podatki $a=3~cm$, $b=4~cm$ in $c=5~cm$. S kosinusnim izrekom izračunajte velikost kota $\gamma$, ki leži nasproti stranice $c$.
        


    \begin{solution}[10cm]

    \end{solution}





    % \bonusquestion[3]
    % \textit{Dodatna naloga:} Pokažite, da vektorja $\overrightarrow{a}=\left(\frac{3}{5},\frac{4}{5}\right)$ in $\overrightarrow{b}=\left(-\frac{4}{5},\frac{3}{5}\right)$ določata ortonormirano bazo ravnine.

    % \begin{solution}[1cm]

    % \end{solution}

\end{questions}



\end{document}